\section{The icosahedron: an exact global certificate}
Put $\varphi=(1+\sqrt5)/2$ and let $I$ be the convex hull of
\[
 (0,\pm1,\pm\varphi),\qquad(\pm1,\pm\varphi,0),\qquad
 (\pm\varphi,0,\pm1).
\]
Its edge length is $2$. A largest section is $I\cap\{z=0\}$. In cyclic
order its vertices are
\[
 (\varphi,0,0),\ (1,\varphi,0),\ (-1,\varphi,0),\
 (-\varphi,0,0),\ (-1,-\varphi,0),\ (1,-\varphi,0).
\]
Thus it is a centrally symmetric, nonregular hexagon. Two opposite sides
have length $2$, and the other four sides have length $\sqrt3$. Its area is
\[
 M=2\varphi(\varphi+1)=2\varphi^3=4+2\sqrt5.
\]
For edge length $s$, the answer is therefore
$\frac{2+\sqrt5}{2}s^2$.

\subsection{A fundamental cone without an assumption about symmetric optima}
The reflections in $x=0$, $y=0$, and
$z=\varphi^2x+\varphi y$ all permute the twelve vertices. For the last one,
use the normal $m=(-\varphi^2,-\varphi,1)$, with
$m\cdot m=4\varphi^2$, and the formula
$R_m(v)=v-2(v\cdot m)m/(m\cdot m)$.
The group they generate is finite: its action on the twelve vertices is
faithful, since those vertices span space. For any nonzero normal $n$,
choose from its finite orbit a vector maximizing its scalar product with
$h=(1,1,10)$. Since $h$ has positive scalar product with each inward
wall normal, reflection in a violated wall would increase that scalar
product. Consequently the chosen vector lies in
\[
 C=\{(x,y,z):x\ge0,\ y\ge0,\ z\ge\varphi^2x+\varphi y\}.
\]
This is a finite comparison argument. It does not assert that a maximum
must itself have a symmetry.

The signs of scalar products with the twelve vertices are constant in
$C_0=C\cap\{x\le\varphi y\}$ and
$C_1=C\cap\{x\ge\varphi y\}$, apart from boundary zeros. Explicit
nonnegative ray coordinates are
\[
\begin{array}{c|ccc}
 &x&y&z\\\hline
C_0&\varphi c&b+c&a+\varphi b+(3\varphi+1)c\\
C_1&b+\varphi c&c&a+\varphi^2b+(3\varphi+1)c
\end{array}
\qquad a,b,c\ge0.
\]
Indeed $a=z-\varphi^2x-\varphi y$ in both cases; in $C_0$ take
$c=x/\varphi$, $b=y-x/\varphi$, and in $C_1$ take $c=y$,
$b=x-\varphi y$. Every nonzero vector in either cone has $z>0$.

\subsection{Area formulas}
For a plane $xx_1+yx_2+zx_3=0$, write
$Q=x^2+y^2+z^2$. Orthogonal projection of its section onto $x_3=0$
multiplies area by $z/\sqrt Q$. On $C_j$ its projected area is $N_j/D_j$,
where
\begin{align*}
N_0={}&2\varphi^4(\varphi z^2+2\varphi yz-x^2)
 \bigl(\varphi^2z^2+(4\varphi-2)yz+(2-\varphi)y^2-x^2\bigr),\\
D_0={}&\bigl((\varphi z+\varphi^2y)^2-x^2\bigr)
        \bigl((\varphi z+(\varphi-1)y)^2-x^2\bigr),\\
N_1={}&2\varphi^4\bigl[x^5+(5-3\varphi)x^4z
 +(\varphi-2)x^3y^2-(2\varphi+1)x^3z^2\\
&\quad +(3\varphi-9)x^2y^2z+(\varphi-2)x^2z^3
 -(3\varphi+9)xy^2z^2\\
&\quad +(2\varphi+1)xz^4+(\varphi+1)y^4z-y^2z^3\bigr],\\
D_1={}&x\bigl((\varphi z-x)^2-\varphi^4y^2\bigr)
        \bigl((\varphi z+x)^2-(\varphi-1)^2y^2\bigr).
\end{align*}
Thus the actual area is $A=\sqrt Q\,N_j/(zD_j)$.

Here is a reproducible derivation, rather than an invocation of an
optimizer. Enumerate pairs of original vertices at squared distance $4$.
For each crossing edge $[u,v]$, its section point is
\[
 q_{uv}=\frac{(n\cdot u)v-(n\cdot v)u}{n\cdot(u-v)}.
\]
Order these points by the adjacency of the original triangular faces,
and apply the planar shoelace formula to their first two coordinates.
The interiors of $C_0$ and $C_1$ have respectively ten and twelve
section vertices. Canceling the shoelace rational functions gives the
four displayed polynomials. The accompanying standard-library verifier
checks these identities by polynomial cross multiplication. It also
checks them directly on all seven relatively open faces of each ray
cone: original vertices on the plane are inserted directly and edges
with endpoints of opposite sign are intersected by the formula above.
Thus boundary cases do not rely on a floating-point perturbation or on
a continuity claim.

For nonzero $n\in C$, $D_0>0$. To see this, factor each quadratic in
$D_0$ into its two linear factors and use $\varphi z>x$.
For $n\in C_1$ with $x>0$, $D_1>0$: its smallest potentially problematic
factor satisfies
\[
 \varphi z-x-\varphi^2y=\varphi a+2\varphi x>0,
\]
and
\[
 \varphi z+x-(\varphi-1)y
 =\varphi a+2\varphi^2x+2y>0.
\]
If $x=0$ in $C_1$, then $y=0$, and the section is the displayed maximizing
hexagon.

\subsection{The global inequality and equality}
Substitute the ray coordinates into
\[
 P_j=M^2 z^2D_j^2-QN_j^2.
\]
These are homogeneous polynomials in $a,b,c$. Their exact expansions
have the following properties:
\[
\begin{array}{c|c|c|c|c}
 j&\deg P_j&\text{nonzero monomials}
 &\min A_\alpha&\min B_\alpha\\\hline
0&10&65&984&440\\
1&12&85&304&136
\end{array}
\]
Every coefficient is $A_\alpha+B_\alpha\sqrt5$, with integer
$A_\alpha,B_\alpha$ at least the corresponding positive minima.
The complete exact coefficient lists accompany this proof as JSON and
CSV; the verifier expands the displayed $N_j,D_j$ and checks equality
with those lists, using only rational arithmetic and $\sqrt5^2=5$.
In particular this is a finite polynomial identity plus coefficientwise
positivity, not a numerical bound. For easy independent strictness
checks, the coefficients of $b^d,c^d$ are
\[
\begin{array}{c|r|r}
j&[b^d]P_j&[c^d]P_j\\\hline
0&1357024+606880\sqrt5&12430643200+5559152640\sqrt5\\
1&307968+137728\sqrt5&7254425600+3244277760\sqrt5.
\end{array}
\]
It follows term by term that $P_j\ge0$ for $a,b,c\ge0$, and that
$P_j>0$ whenever $b+c>0$. Dividing by the positive quantity $z^2D_j^2$
gives $A^2\le M^2$, hence $A\le M$. Equality in this fundamental cone
requires $b=c=0$, so the normal is the $z$-axis. Every central maximizing
section is therefore congruent to the displayed hexagon. Together with
the central-section reduction for centrally symmetric polytopes, this
establishes the bound for every affine section.

\subsection{Arithmetic strength}
All inputs, matrices, vertices, and certificates lie in
$\mathbb Q(\sqrt5)$. Represent an element by two rational numbers; represent
a rational by a signed numerator and a positive denominator. Comparison
is decidable by rational comparisons and, when signs conflict, comparison
of two integer squares. Polynomial evaluation and expansion are bounded
loops over fixed finite lists. The reflection group can be represented
by permutations of twelve vertices; closure and finite maximization
are bounded by $12!$. The section graph is constructed by bounded
search through $30$ edges and $20$ faces. Every loop bound and every
arithmetic operation is primitive recursive, and its correctness uses
quantifier-free induction. Positive coefficient evaluation is a fixed
finite sum of nonnegative products. Thus the displayed global inequality
has a direct iPRA derivation for coded algebraic planes; the common
geometric interpretation and approximation lemmas handle the chosen
constructive interpretation of arbitrary real planes. Lean should first
formalize the finite $\mathbb Q(\sqrt5)$ certificate checker, prove its
soundness in the arithmetized iPRA syntax, and then instantiate it with
the included coefficient data. Running Python verifies the mathematical
certificate; it is not itself a completed Lean kernel check or an
object-language iPRA proof term.
